% Chapter 4 (Design) of the essay. Compile from the infor_sys folder with:
%   tectonic -X compile tieuluan/chapter4/design_04.tex --outdir tieuluan/build
% Every fact about classes, tables and packages comes from tieuluan/tools/model.py;
% figures come from the Visual Paradigm project, numbers from files in tieuluan/build.
\documentclass[12pt,a4paper]{article}
\usepackage{fontspec}
\IfFontExistsTF{Times New Roman}{\setmainfont{Times New Roman}}{\setmainfont{TeX Gyre Termes}}
\IfFontExistsTF{Arial}{\setsansfont{Arial}}{\setsansfont{TeX Gyre Heros}}
\IfFontExistsTF{Menlo}{\setmonofont{Menlo}[Scale=0.80]}{\setmonofont{TeX Gyre Cursor}[Scale=0.85]}
\usepackage[left=2.6cm,right=1.8cm,top=2.4cm,bottom=2.4cm,headheight=16pt]{geometry}
\usepackage{setspace,titlesec,tikz}
\setstretch{1.2}
\usepackage{graphicx,array,booktabs,longtable,tabularx,multirow}
\usepackage{amsmath,amssymb}
\usepackage{enumitem,float,caption,needspace,pdflscape}
\usepackage{listings,fancyvrb}
\usepackage{xcolor}
\definecolor{ptitred}{HTML}{A51C30}
\definecolor{codegray}{HTML}{F6F6F6}
\usepackage[obeyspaces,spaces]{url}
\usepackage[unicode,hidelinks]{hyperref}
\usepackage{fancyhdr}
\pagestyle{fancy}
\fancyhf{}
\fancyhead[L]{\small\itshape Tiểu luận Phân tích và thiết kế hệ thống thông tin}
\fancyhead[R]{\small\itshape Chương 4. Thiết kế (Design)}
\fancyfoot[C]{\thepage}
\fancypagestyle{plain}{\fancyhf{}\fancyhead[L]{\small\itshape Tiểu luận Phân tích và thiết kế hệ thống thông tin}%
  \fancyhead[R]{\small\itshape Chương 4. Thiết kế (Design)}\fancyfoot[C]{\thepage}}
\renewcommand{\contentsname}{Mục lục}
\renewcommand{\figurename}{Hình}
\renewcommand{\tablename}{Bảng}
\renewcommand{\listfigurename}{Danh mục hình vẽ}
\renewcommand{\listtablename}{Danh mục bảng}
\renewcommand{\thesection}{4.\arabic{section}}
\renewcommand{\thefigure}{4.\arabic{figure}}
\renewcommand{\thetable}{4.\arabic{table}}
\setlength{\parindent}{0.8cm}
\setlength{\parskip}{4pt}
% Numbered headings in the usual Vietnamese style: "4.1." and "4.1.1." in black bold.
\titleformat{\section}{\normalfont\large\bfseries}{\thesection.}{0.5em}{}
\titleformat{\subsection}{\normalfont\normalsize\bfseries}{\thesubsection.}{0.5em}{}
\titlespacing*{\section}{0pt}{14pt}{6pt}
\titlespacing*{\subsection}{0pt}{10pt}{4pt}
\setlength{\emergencystretch}{3em}
\renewcommand{\arraystretch}{1.18}
\setlength{\tabcolsep}{4.6pt}
\newcolumntype{P}[1]{>{\raggedright\arraybackslash}p{#1}}
\renewcommand{\tabularxcolumn}[1]{>{\raggedright\arraybackslash}p{#1}}
\setlist{topsep=3pt,itemsep=2pt,parsep=0pt}
\captionsetup{font=small,labelfont=bf,labelsep=period}
\captionsetup[figure]{textfont=it}
\setlength{\LTpre}{4pt}
\setlength{\LTpost}{6pt}

% Identifiers and signatures: monospaced and breakable after . , ( ) : < > _
\makeatletter
\def\UrlBreaks{\do\.\do\,\do\(\do\)\do\:\do\<\do\>\do\_\do\/\do\[\do\]}
\makeatother
\DeclareUrlCommand\code{\urlstyle{tt}}

\lstdefinestyle{ecom}{basicstyle=\ttfamily\scriptsize,breaklines=true,frame=single,
  backgroundcolor=\color{codegray},columns=fullflexible,keepspaces=true,showstringspaces=false,
  numbers=left,numberstyle=\tiny\color{gray},xleftmargin=2em,framexleftmargin=1.6em,tabsize=4,
  keywordstyle=\bfseries,commentstyle=\itshape\color{black!65}}
\lstset{style=ecom}
\fvset{fontsize=\scriptsize,frame=single,framesep=6pt}

% Figure exported from the Visual Paradigm project; a placeholder keeps drafts compiling.
\newcommand{\vpimage}[2]{%
  \IfFileExists{../figures/vp/#1}%
    {\includegraphics[width=\linewidth,height=#2,keepaspectratio]{../figures/vp/#1}}%
    {\fbox{\parbox[c][5cm][c]{0.85\linewidth}{\centering Hình đang được dựng từ Visual Paradigm\\[3pt]
      \small\texttt{\detokenize{#1}}}}}}
\newcommand{\txtfile}[1]{\IfFileExists{#1}{\VerbatimInput{#1}}{\fbox{Chưa có tệp \texttt{\detokenize{#1}}}}}
\newcommand{\codefile}[2]{\IfFileExists{#2}{\lstinputlisting[language=#1]{#2}}{\fbox{Chưa có tệp \texttt{\detokenize{#2}}}}}

% Code and SQL statistics are written by tools/gen_java.py and tools/gen_sql.py.
\IfFileExists{../build/code_stats.tex}{\newcommand{\javaFileCount}{44}
\newcommand{\javaLineCount}{1719}
\newcommand{\javaPackageCount}{11}
\newcommand{\sqlTableCount}{20}
\newcommand{\sqlLineCount}{200}
}{}
\providecommand{\javaFileCount}{?}
\providecommand{\javaLineCount}{?}
\providecommand{\javaPackageCount}{?}
\providecommand{\sqlTableCount}{?}
\providecommand{\sqlLineCount}{?}

\begin{document}
\begin{titlepage}
\begin{tikzpicture}[remember picture,overlay]
  \draw[line width=1.6pt] ([xshift=1.6cm,yshift=-1.4cm]current page.north west)
    rectangle ([xshift=-1.2cm,yshift=1.4cm]current page.south east);
  \draw[line width=0.5pt] ([xshift=1.75cm,yshift=-1.55cm]current page.north west)
    rectangle ([xshift=-1.35cm,yshift=1.55cm]current page.south east);
\end{tikzpicture}
\centering
\setstretch{1.15}
{\large\bfseries HỌC VIỆN CÔNG NGHỆ BƯU CHÍNH VIỄN THÔNG\par}
\vspace{2pt}\rule{0.35\linewidth}{0.6pt}\par
\vspace{2.6cm}
{\Large\bfseries TIỂU LUẬN MÔN HỌC\par}
\vspace{0.25cm}
{\Large\bfseries PHÂN TÍCH VÀ THIẾT KẾ HỆ THỐNG THÔNG TIN\par}
\vspace{0.9cm}
\rule{0.55\linewidth}{0.6pt}\par
\vspace{0.9cm}
{\large\bfseries ĐỀ TÀI\par}
\vspace{0.4cm}
{\LARGE\bfseries TỪ YÊU CẦU ĐẾN LẬP TRÌNH\par}
\vspace{0.3cm}
{\large\itshape Phân tích và thiết kế ứng dụng thương mại điện tử\\ có tìm kiếm đa phương thức\par}
\vspace{1.4cm}
{\Large\bfseries CHƯƠNG 4. THIẾT KẾ HỆ THỐNG (DESIGN)\par}
\vspace{2.4cm}
{\large
\begin{tabular}{@{}ll@{}}
Sinh viên thực hiện: & Vũ Hùng Anh\\[0.25cm]
Mã sinh viên: & B23DCDT022\\[0.25cm]
Lớp: & 04
\end{tabular}\par}
\vfill
{\large Hà Nội, tháng 10 năm 2026\par}
\vspace{0.6cm}
\end{titlepage}

{\setlength{\parskip}{0pt}\setstretch{1.1}\tableofcontents}
\clearpage
{\setlength{\parskip}{0pt}\setstretch{1.1}\listoffigures\vspace{0.6cm}\listoftables}
\clearpage

%% ===================================================================== intro
\begin{center}{\Large\bfseries CHƯƠNG 4. THIẾT KẾ HỆ THỐNG (DESIGN)}\end{center}
\addcontentsline{toc}{section}{Chương 4. Thiết kế hệ thống (Design)}

Chương này được viết lại theo dàn ý 4.1--4.6 mà giảng viên giao cho phần Design, thay cho chương
Design của bản tiểu luận V1. Bản V1 trình bày nguyên tắc thiết kế bằng lời và chưa có sơ đồ thiết kế;
bản này đi qua một chuỗi sản phẩm cụ thể: sơ đồ lớp phân tích được cập nhật, mô hình dữ liệu và cơ sở
dữ liệu sinh từ nó, sơ đồ lớp thiết kế, kiến trúc phân tầng, và mã nguồn Java sinh từ mô hình phân tầng.

Nguồn đầu vào gồm: mô hình 19 lớp của Bài kiểm tra 01 (ba project Visual Paradigm
\code{ecomAnalysis}, \code{eComDataModel}, \code{eComDesign}); slide A03 \emph{Domain Model and
Analysis Class Diagram} (Cart, CartItem, Order, OrderItem và phép biến đổi checkout); slide A05
\emph{From Analysis to Design} (Controller -- Service -- Domain -- Repository, layer và package);
và prototype tìm kiếm đa phương thức Assignment 06. Toàn bộ sơ đồ được dựng trong Visual Paradigm 18.1
Community Edition qua Open API thành một project chỉnh sửa được; hình trong bài là ảnh xuất trực tiếp
từ project đó. Lớp, thuộc tính, phương thức, bảng và khóa đều lấy từ một tệp đặc tả chung
(\code{tools/model.py}), nên sơ đồ, câu lệnh SQL và mã Java dùng đúng cùng một tên gọi.

\begin{table}[H]
\centering\small
\caption{Luồng sản phẩm của chương Design}\label{tab:flow}
\begin{tabular}{@{}P{1.0cm}P{3.6cm}P{4.4cm}P{6.2cm}@{}}
\toprule
Mục & Câu hỏi trả lời & Đầu vào & Sản phẩm\\
\midrule
4.1 & Hệ thống cần những khái niệm nào? & 19 lớp Bài kiểm tra 01, slide A03 & Sơ đồ lớp phân tích 20 lớp (Hình~\ref{fig:analysis})\\
4.2 & Dữ liệu được lưu như thế nào? & Sơ đồ 4.1 & Data model ORM (Hình~\ref{fig:orm}), mô hình CSDL 20 bảng (Hình~\ref{fig:erd}), tệp \code{db/ecom_schema.sql}\\
4.3 & Mỗi lớp giữ dữ liệu và hành vi gì? & Sơ đồ 4.1, slide A05 & Sơ đồ lớp thiết kế có kiểu, phương thức, điều hướng (Hình~\ref{fig:design})\\
4.4 & Phần mềm chia thành tầng nào, lớp nào gọi lớp nào? & Sơ đồ 4.3, prototype A6 & Sơ đồ layer (Hình~\ref{fig:layers}), sơ đồ lớp theo tầng (Hình~\ref{fig:layerclasses}), sequence checkout (Hình~\ref{fig:seq})\\
4.5 & Mã nguồn tương ứng ra sao? & Mô hình 4.4 & Mã Java trong \code{code/src}, chia theo package của 4.4\\
\bottomrule
\end{tabular}
\end{table}

%% ===================================================================== 4.1
\section{Sơ đồ lớp phân tích (Class Diagram in Analysis)}\label{sec:41}

Sơ đồ lớp phân tích trả lời câu hỏi \emph{hệ thống cần biểu diễn những khái niệm nào và chúng liên hệ
với nhau ra sao}. Theo slide A03, ở mức này chưa đưa vào Controller, Service, Repository hay bảng dữ
liệu; mỗi lớp là một khái niệm nghiệp vụ, mỗi quan hệ là một ràng buộc nghiệp vụ có bội số.

\subsection{Cập nhật sơ đồ so với Bài kiểm tra 01 (Update)}

Sơ đồ của Bài kiểm tra 01 dùng đúng 19 lớp do đề cho. Khi dùng nó làm đầu vào cho thiết kế, có một số
điểm làm phần sau phức tạp không cần thiết, nên sơ đồ được cập nhật như Bảng~\ref{tab:changes}. Kết quả
có 20 lớp: 19 lớp cũ và lớp trừu tượng \code{Product}.

\begin{longtable}{@{}P{2.6cm}P{3.8cm}P{4.4cm}P{4.7cm}@{}}
\caption{Các thay đổi của sơ đồ phân tích so với Bài kiểm tra 01}\label{tab:changes}\\
\toprule
Điểm thay đổi & Bài kiểm tra 01 & Bản cập nhật & Lý do\\
\midrule
\endfirsthead
\toprule
Điểm thay đổi & Bài kiểm tra 01 & Bản cập nhật & Lý do\\
\midrule
\endhead
\small Lớp cha của hàng hóa & \small Không có \code{Product}; \code{CartItem} và \code{OrderItem} mỗi lớp có bốn liên kết tới \code{Book}, \code{Clothes}, \code{Shoes}, \code{Electronics} với ràng buộc \{xor\} &
\small Thêm lớp trừu tượng \code{Product}, là lớp cha của \code{Book}, \code{Clothes}, \code{Shoes}, \code{Electronics}. \code{CartItem} và \code{OrderItem} chỉ còn một liên kết \emph{refers to} tới \code{Product}, bội số 0..* -- 1 &
\small Một dòng hàng luôn chỉ tới đúng một sản phẩm mà không cần ràng buộc \{xor\}; khớp mô hình Cart $\to$ CartItem $\to$ Product của slide A03; khi lưu trữ chỉ cần một khóa ngoại \code{product_id}\\
\small Thuộc tính chung của hàng hóa & \small \code{productId}, \code{price}, \code{stockQuantity} và tên hàng lặp lại ở bốn lớp; \code{Book} dùng \code{title} &
\small Chuyển \code{productId}, \code{name}, \code{price}, \code{stockQuantity} lên \code{Product}; \code{Book.title} thành \code{Product.name}; thêm \code{Product.imageUrl} &
\small Bỏ trùng lặp; \code{imageUrl} cần cho chức năng tìm kiếm bằng ảnh ở Chương 3\\
\small Thẻ thanh toán & \small \code{PayCredit} có \code{cardNumber}, \code{cardHolder}, \code{bankName}, \code{expiryDate} &
\small \code{PayCredit} có \code{cardHolder}, \code{bankName}, \code{cardLast4}, \code{authorizationCode} &
\small Theo quyết định D-01 (Chương 1), hệ thống không lưu số thẻ đầy đủ; chỉ giữ bốn số cuối và mã cấp phép do cổng thanh toán trả về\\
\small Bội số và tên quan hệ & \small Chỉ bổ sung ở data model (Câu 2) &
\small Ghi ngay trên sơ đồ phân tích; \code{Order} -- \code{Payment} là 1 -- 0..* &
\small Slide A03 yêu cầu ghi bội số ở mức phân tích; một đơn có thể có nhiều lần thanh toán khi lần trước thất bại\\
\small Phụ thuộc checkout & \small \code{Cart} $\dashrightarrow$ \code{Order} «create» \emph{checkout} &
\small Giữ nguyên &
\small Slide A03: checkout là một phép biến đổi Cart thành Order, không phải liên kết cố định\\
\bottomrule
\end{longtable}

\subsection{Mô tả sơ đồ bằng lời (Describe in words)}

\begin{landscape}
\begin{figure}[p]
\centering
\vpimage{c41_analysis.png}{0.80\textheight}
\caption[Sơ đồ lớp phân tích (Analysis Class Diagram)]{Sơ đồ lớp phân tích đã cập nhật (20 lớp, Visual Paradigm). Lớp in nghiêng là lớp trừu tượng.}\label{fig:analysis}
\end{figure}
\end{landscape}

Hình~\ref{fig:analysis} gồm bốn nhóm lớp.

\textbf{Nhóm khách hàng.} \code{Customer} là người mua, định danh bằng \code{customerId} và có
\code{email}, \code{phone}, \code{dateOfBirth}. Họ tên được tách thành đối tượng giá trị
\code{FullName} (\code{firstName}, \code{middleName}, \code{lastName}). Mỗi khách có một hoặc nhiều
\code{Address} (\code{houseNumber}, \code{street}, \code{district}, \code{city}). Hai hạng khách là
\code{CustomerVIP} (\code{vipLevel}, \code{discountRate}, \code{rewardPoints}) và \code{CustomerNew}
(\code{registeredDate}, \code{welcomeVoucher}).

\textbf{Nhóm sản phẩm.} \code{Product} là lớp trừu tượng chứa phần chung của mọi mặt hàng:
\code{productId}, \code{name}, \code{price}, \code{stockQuantity}, \code{imageUrl}. Bốn nhóm hàng
chuyên biệt hóa \code{Product}: \code{Book} (\code{author}, \code{publisher}), \code{Clothes}
(\code{brand}, \code{size}, \code{color}, \code{material}), \code{Shoes} (\code{brand}, \code{size},
\code{color}) và \code{Electronics} (\code{brand}, \code{warrantyMonths}). \code{Laptop} (\code{cpu},
\code{ram}, \code{storage}, \code{screenSize}) và \code{Mobile} (\code{os}, \code{camera},
\code{screenSize}, \code{batteryCapacity}) chuyên biệt hóa tiếp \code{Electronics}.

\textbf{Nhóm giỏ hàng và đơn hàng.} \code{Cart} (\code{cartId}, \code{createdDate},
\code{updatedDate}) là trạng thái mua sắm tạm thời, gồm các \code{CartItem} (\code{quantity},
\code{cartPrice}). \code{Order} (\code{orderId}, \code{orderDate}, \code{status}, \code{totalAmount}) là
giao dịch đã xác nhận, gồm các \code{OrderItem} (\code{quantity}, \code{unitPrice}).

\textbf{Nhóm giao hàng và thanh toán.} \code{Shipping} (\code{shippingId}, \code{carrier},
\code{trackingNumber}, \code{shippingFee}, \code{status}, \code{deliveredDate}) ghi thông tin giao một
đơn tới một địa chỉ. \code{Payment} là lớp trừu tượng (\code{paymentId}, \code{amount},
\code{paymentDate}, \code{status}) với hai hình thức \code{PayCash} (\code{cashReceived},
\code{changeAmount}) và \code{PayCredit} (\code{cardHolder}, \code{bankName}, \code{cardLast4},
\code{authorizationCode}).

\begin{longtable}{@{}P{2.6cm}P{6.5cm}P{6.8cm}@{}}
\caption{Ý nghĩa của 20 lớp phân tích}\label{tab:classes}\\
\toprule
Lớp & Ý nghĩa nghiệp vụ & Thuộc tính\\
\midrule
\endfirsthead
\toprule
Lớp & Ý nghĩa nghiệp vụ & Thuộc tính\\
\midrule
\endhead
\small\code{Customer} & \small Người mua có tài khoản & \small\code{customerId}, \code{email}, \code{phone}, \code{dateOfBirth}\\
\small\code{FullName} & \small Họ tên của khách, đối tượng giá trị & \small\code{firstName}, \code{middleName}, \code{lastName}\\
\small\code{Address} & \small Địa chỉ liên hệ và nhận hàng & \small\code{houseNumber}, \code{street}, \code{district}, \code{city}\\
\small\code{CustomerVIP} & \small Khách hạng VIP, được chiết khấu và tích điểm & \small\code{vipLevel}, \code{discountRate}, \code{rewardPoints}\\
\small\code{CustomerNew} & \small Khách mới đăng ký, có thể có voucher chào mừng & \small\code{registeredDate}, \code{welcomeVoucher}\\
\small\code{Product} & \small Mặt hàng được bán (trừu tượng) & \small\code{productId}, \code{name}, \code{price}, \code{stockQuantity}, \code{imageUrl}\\
\small\code{Book} & \small Sách & \small\code{author}, \code{publisher}\\
\small\code{Clothes} & \small Quần áo & \small\code{brand}, \code{size}, \code{color}, \code{material}\\
\small\code{Shoes} & \small Giày & \small\code{brand}, \code{size}, \code{color}\\
\small\code{Electronics} & \small Hàng điện tử & \small\code{brand}, \code{warrantyMonths}\\
\small\code{Laptop} & \small Máy tính xách tay & \small\code{cpu}, \code{ram}, \code{storage}, \code{screenSize}\\
\small\code{Mobile} & \small Điện thoại di động & \small\code{os}, \code{camera}, \code{screenSize}, \code{batteryCapacity}\\
\small\code{Cart} & \small Giỏ hàng đang mua dở & \small\code{cartId}, \code{createdDate}, \code{updatedDate}\\
\small\code{CartItem} & \small Một sản phẩm và số lượng định mua trong giỏ & \small\code{quantity}, \code{cartPrice} (giá lúc thêm vào giỏ)\\
\small\code{Order} & \small Đơn hàng đã xác nhận & \small\code{orderId}, \code{orderDate}, \code{status}, \code{totalAmount}\\
\small\code{OrderItem} & \small Một sản phẩm đã mua trong đơn & \small\code{quantity}, \code{unitPrice} (giá chốt lúc đặt)\\
\small\code{Shipping} & \small Thông tin giao hàng của một đơn & \small\code{shippingId}, \code{carrier}, \code{trackingNumber}, \code{shippingFee}, \code{status}, \code{deliveredDate}\\
\small\code{Payment} & \small Một lần thanh toán (trừu tượng) & \small\code{paymentId}, \code{amount}, \code{paymentDate}, \code{status}\\
\small\code{PayCash} & \small Thanh toán tiền mặt & \small\code{cashReceived}, \code{changeAmount}\\
\small\code{PayCredit} & \small Thanh toán thẻ qua cổng thanh toán & \small\code{cardHolder}, \code{bankName}, \code{cardLast4}, \code{authorizationCode}\\
\bottomrule
\end{longtable}

\subsection{Giải thích các quan hệ (Explain relationships)}

Bảng~\ref{tab:assoc41} giải thích từng quan hệ trên Hình~\ref{fig:analysis}. Bội số ghi ở một đầu cho
biết \emph{một} đối tượng ở đầu đối diện liên kết với bao nhiêu đối tượng ở đầu này.

\begin{longtable}{@{}P{4.0cm}P{2.3cm}P{1.8cm}P{7.4cm}@{}}
\caption{Quan hệ giữa các lớp phân tích}\label{tab:assoc41}\\
\toprule
Quan hệ (tên) & Loại & Bội số & Ý nghĩa nghiệp vụ\\
\midrule
\endfirsthead
\toprule
Quan hệ (tên) & Loại & Bội số & Ý nghĩa nghiệp vụ\\
\midrule
\endhead
\small\code{Customer} -- \code{FullName} (has) & \small Composition & \small 1 -- 1 & \small Mỗi khách có đúng một họ tên; họ tên không tồn tại ngoài khách hàng nên mất cùng khách.\\
\small\code{Customer} -- \code{Address} (has) & \small Aggregation & \small 1 -- 1..* & \small Khách phải có ít nhất một địa chỉ để nhận hàng. Địa chỉ còn được \code{Shipping} tham chiếu nên không bị xóa theo khách, vì vậy dùng aggregation chứ không dùng composition.\\
\small\code{Customer} -- \code{Cart} (owns) & \small Association & \small 1 -- 0..1 & \small Mỗi khách có tối đa một giỏ đang hoạt động; khách mới có thể chưa có giỏ.\\
\small\code{Customer} -- \code{Order} (places) & \small Association & \small 1 -- 0..* & \small Khách đặt nhiều đơn; khách mới chưa có đơn. Đơn tồn tại độc lập để còn đối soát.\\
\small\code{Cart} -- \code{CartItem} (contains) & \small Composition & \small 1 -- 0..* & \small Dòng giỏ chỉ tồn tại trong giỏ; giỏ rỗng là trạng thái hợp lệ.\\
\small\code{Order} -- \code{OrderItem} (contains) & \small Composition & \small 1 -- 1..* & \small Dòng đơn là bộ phận của đơn; đơn không có dòng hàng nào là vô nghĩa.\\
\small\code{CartItem} -- \code{Product} (refers to) & \small Association & \small 0..* -- 1 & \small Mỗi dòng giỏ chỉ tới đúng một sản phẩm; một sản phẩm có thể nằm trong nhiều giỏ.\\
\small\code{OrderItem} -- \code{Product} (refers to) & \small Association & \small 0..* -- 1 & \small Mỗi dòng đơn chỉ tới đúng một sản phẩm; một sản phẩm có thể có trong nhiều đơn.\\
\small\code{Order} -- \code{Shipping} (is delivered by) & \small Association & \small 1 -- 0..1 & \small Đơn mới tạo có thể chưa có thông tin giao; khi giao thì có đúng một bản ghi giao hàng.\\
\small\code{Shipping} -- \code{Address} (ships to) & \small Association & \small 0..* -- 1 & \small Mỗi lần giao tới đúng một địa chỉ; một địa chỉ nhận được nhiều lần giao.\\
\small\code{Order} -- \code{Payment} (is paid by) & \small Association & \small 1 -- 0..* & \small Đơn có thể chưa thanh toán, hoặc có nhiều lần thanh toán khi lần trước thất bại.\\
\small\code{CustomerVIP}, \code{CustomerNew} $\to$ \code{Customer} & \small Generalization & & \small Hai hạng khách dùng lại thuộc tính của \code{Customer} và thêm thuộc tính riêng.\\
\small\code{Book}, \code{Clothes}, \code{Shoes}, \code{Electronics} $\to$ \code{Product} & \small Generalization & & \small Bốn nhóm hàng dùng chung mã, tên, giá, tồn kho, ảnh; mỗi nhóm thêm thuộc tính riêng.\\
\small\code{Laptop}, \code{Mobile} $\to$ \code{Electronics} & \small Generalization & & \small Hai loại thiết bị kế thừa \code{brand}, \code{warrantyMonths} và phần chung của \code{Product}.\\
\small\code{PayCash}, \code{PayCredit} $\to$ \code{Payment} & \small Generalization & & \small Hai hình thức thanh toán dùng chung số tiền, thời điểm, trạng thái.\\
\small\code{Cart} $\dashrightarrow$ \code{Order} (checkout) & \small Dependency «create» & & \small Checkout dùng nội dung giỏ để tạo một đơn mới; sau đó giỏ được làm rỗng.\\
\bottomrule
\end{longtable}

\code{Product} và \code{Payment} là lớp trừu tượng (tên in nghiêng) vì không có mặt hàng nào không thuộc
một nhóm cụ thể và không có lần thanh toán nào không thuộc một hình thức cụ thể.

\textbf{CartItem khác OrderItem, và checkout là một phép biến đổi.} Hai lớp có thuộc tính giống nhau
nhưng khác ý nghĩa (slide A03). \code{CartItem} là thông tin mua sắm tạm thời: \code{cartPrice} là giá
lúc khách thêm hàng vào giỏ và có thể thay đổi. \code{OrderItem} là thông tin giao dịch lịch sử:
\code{unitPrice} là giá được chốt khi tạo đơn và không đổi khi giá catalog thay đổi sau đó. Checkout
biến giỏ $\{\text{CartItem}_1,\dots,\text{CartItem}_n\}$ thành đơn
$\{\text{OrderItem}_1,\dots,\text{OrderItem}_n\}$ với cùng sản phẩm và số lượng, sau đó giỏ trở thành
rỗng. Vì giỏ và đơn là hai trạng thái nghiệp vụ khác nhau, quan hệ giữa \code{Cart} và \code{Order} được
vẽ là phụ thuộc «create» thay vì một association cố định.

%% ===================================================================== 4.2
\section{Mô hình dữ liệu và cơ sở dữ liệu (Data Model -- Database)}\label{sec:42}

Mục này chuyển sơ đồ phân tích thành mô hình dữ liệu (data model) rồi sinh cơ sở dữ liệu. Có ba sản
phẩm: data model dạng lớp ORM trong Visual Paradigm (Hình~\ref{fig:orm}), mô hình cơ sở dữ liệu vật
lý gồm 20 bảng (Hình~\ref{fig:erd}) và tệp DDL \code{db/ecom_schema.sql}.

\subsection{Mô hình dữ liệu ORM (Data model)}

\begin{landscape}
\begin{figure}[p]
\centering
\vpimage{c42_orm.png}{0.80\textheight}
\caption[Mô hình dữ liệu ORM-Persistable (Data model)]{Data model: 20 lớp ORM-Persistable, quan hệ kế thừa đã chuyển thành association \emph{is a} 1 -- 0..1.}\label{fig:orm}
\end{figure}
\end{landscape}

Data model trong Hình~\ref{fig:orm} được tạo từ Hình~\ref{fig:analysis} theo cùng bốn bước như Câu 2
của Bài kiểm tra 01:
\begin{enumerate}
\item mọi lớp được gắn stereotype \code{<<ORM-Persistable>>}, nghĩa là đối tượng của lớp được lưu thành
bản ghi;
\item mọi association có tên và bội số ở hai đầu;
\item mỗi quan hệ kế thừa được đổi thành association \emph{is a}, đầu lớp cha có bội số 1, đầu lớp con
có bội số 0..1, vì mô hình quan hệ không có khái niệm kế thừa;
\item phụ thuộc «create» \emph{checkout} bị bỏ vì nó mô tả một thao tác, không phải dữ liệu cần lưu.
\end{enumerate}

\subsection{Quy tắc ánh xạ lớp sang bảng (From data model to database)}

\begin{longtable}{@{}P{3.5cm}P{6.6cm}P{5.8cm}@{}}
\caption{Quy tắc chuyển data model thành bảng}\label{tab:mapping}\\
\toprule
Trường hợp & Quy tắc & Ví dụ\\
\midrule
\endfirsthead
\toprule
Trường hợp & Quy tắc & Ví dụ\\
\midrule
\endhead
\small Lớp thực thể & \small Mỗi lớp thành một bảng, tên viết thường nối bằng gạch dưới. \code{Order} đổi thành \code{orders} vì \code{ORDER} là từ khóa SQL. & \small \code{CartItem} $\to$ \code{cart_item}\\
\small Định danh & \small Thuộc tính định danh thành khóa chính kiểu \code{varchar(20)}. & \small \code{customer_id}, \code{product_id}, \code{order_id}\\
\small Đối tượng giá trị 1 -- 1 & \small Composition 1 -- 1 với \code{FullName} được nhúng thành cột của bảng chủ, không tạo bảng riêng. & \small \code{customer.first_name}, \code{middle_name}, \code{last_name}\\
\small Kế thừa & \small Mỗi lớp một bảng (class table inheritance): bảng con dùng khóa chính đồng thời là khóa ngoại tới bảng cha. & \small \code{book.product_id} $\to$ \code{product}; \code{laptop.product_id} $\to$ \code{electronics}\\
\small Cột phân loại & \small Bảng cha có cột cho biết bản ghi thuộc lớp con nào. & \small \code{product.product_type}, \code{payment.method}\\
\small Association 1 -- * & \small Khóa ngoại đặt ở bảng phía nhiều. & \small \code{orders.customer_id}, \code{payment.order_id}\\
\small Association 1 -- 0..1 & \small Khóa ngoại kèm ràng buộc \code{UNIQUE}. & \small \code{cart.customer_id}, \code{shipping.order_id}\\
\small Composition có dòng hàng & \small Bảng dòng hàng dùng khóa chính ghép; xóa bản ghi cha thì xóa luôn dòng hàng (\code{ON DELETE CASCADE}). & \small \code{cart_item(cart_id, product_id)}, \code{order_item(order_id, product_id)}\\
\small Tiền, số lượng, trạng thái & \small Tiền dùng \code{numeric(12,2)}; số lượng là \code{integer} có \code{CHECK}; trạng thái là chuỗi có \code{CHECK ... IN (...)}. & \small \code{order_item.unit_price}, \code{cart_item.quantity > 0}\\
\small Bảng kỹ thuật & \small Bảng \code{product_feature} lưu vector đặc trưng ảnh cho phân hệ tìm kiếm; không có trong mô hình phân tích, được thêm vì thiết kế gói search (mục 4.4.4). & \small \code{product_feature(product_id, model_version, feature_vector, updated_at)}\\
\bottomrule
\end{longtable}

\subsection{Giải thích các bảng (Explain tables)}

\begin{landscape}
\begin{figure}[p]
\centering
\vpimage{c42_tables.png}{0.80\textheight}
\caption[Mô hình cơ sở dữ liệu vật lý gồm 20 bảng]{Mô hình cơ sở dữ liệu vật lý của eCom (Visual Paradigm): 20 bảng «table»; mỗi cột ghi kiểu SQL và \{PK\} khóa chính, \{FK\} khóa ngoại, \{U\} duy nhất, \{NULL\} cho phép rỗng; mỗi đường nối là một khóa ngoại, ghi tên quan hệ, cột khóa ngoại và bội số.}\label{fig:erd}
\end{figure}
\end{landscape}

Hình~\ref{fig:erd} là mô hình cơ sở dữ liệu vật lý dựng trong Visual Paradigm theo các quy tắc ở
Bảng~\ref{tab:mapping}. Mỗi bảng được vẽ như một lớp có stereotype «table»: tên cột, kiểu SQL và các khóa
nằm ngay trong hộp, nên đọc được toàn bộ lược đồ trên một hình. Đường nối đi từ bảng cha tới bảng con,
ghi tên quan hệ và cột khóa ngoại, ví dụ \emph{contains (order\_id)} từ \code{orders} tới
\code{order_item}. Project VP còn có sơ đồ ERD (4.2c) dựng từ cùng đặc tả bằng phần tử bảng, cột và khóa
ngoại của VP. Bảng~\ref{tab:tables} liệt kê 20 bảng cùng lớp nguồn của chúng.

\begin{longtable}{@{}P{0.6cm}P{2.8cm}P{3.2cm}P{5.1cm}P{3.4cm}@{}}
\caption{Danh sách bảng của cơ sở dữ liệu}\label{tab:tables}\\
\toprule
\# & Bảng & Khóa chính & Khóa ngoại $\to$ bảng cha & Lớp nguồn\\
\midrule
\endfirsthead
\toprule
\# & Bảng & Khóa chính & Khóa ngoại $\to$ bảng cha & Lớp nguồn\\
\midrule
\endhead
\small 1 & \small\code{customer} & \small\code{customer_id} & \small -- & \small \code{Customer} + \code{FullName}\\
\small 2 & \small\code{customer_vip} & \small\code{customer_id} & \small\code{customer_id} $\to$ \code{customer} & \small\code{CustomerVIP}\\
\small 3 & \small\code{customer_new} & \small\code{customer_id} & \small\code{customer_id} $\to$ \code{customer} & \small\code{CustomerNew}\\
\small 4 & \small\code{address} & \small\code{address_id} & \small\code{customer_id} $\to$ \code{customer} & \small\code{Address}\\
\small 5 & \small\code{product} & \small\code{product_id} & \small -- & \small\code{Product}\\
\small 6 & \small\code{book} & \small\code{product_id} & \small\code{product_id} $\to$ \code{product} & \small\code{Book}\\
\small 7 & \small\code{clothes} & \small\code{product_id} & \small\code{product_id} $\to$ \code{product} & \small\code{Clothes}\\
\small 8 & \small\code{shoes} & \small\code{product_id} & \small\code{product_id} $\to$ \code{product} & \small\code{Shoes}\\
\small 9 & \small\code{electronics} & \small\code{product_id} & \small\code{product_id} $\to$ \code{product} & \small\code{Electronics}\\
\small 10 & \small\code{laptop} & \small\code{product_id} & \small\code{product_id} $\to$ \code{electronics} & \small\code{Laptop}\\
\small 11 & \small\code{mobile} & \small\code{product_id} & \small\code{product_id} $\to$ \code{electronics} & \small\code{Mobile}\\
\small 12 & \small\code{product_feature} & \small\code{product_id} & \small\code{product_id} $\to$ \code{product} & \small \code{VectorIndex} (gói search)\\
\small 13 & \small\code{cart} & \small\code{cart_id} & \small\code{customer_id} $\to$ \code{customer} (UNIQUE) & \small\code{Cart}\\
\small 14 & \small\code{cart_item} & \small\code{cart_id}, \code{product_id} & \small\code{cart_id} $\to$ \code{cart}; \code{product_id} $\to$ \code{product} & \small\code{CartItem}\\
\small 15 & \small\code{orders} & \small\code{order_id} & \small\code{customer_id} $\to$ \code{customer} & \small\code{Order}\\
\small 16 & \small\code{order_item} & \small\code{order_id}, \code{product_id} & \small\code{order_id} $\to$ \code{orders}; \code{product_id} $\to$ \code{product} & \small\code{OrderItem}\\
\small 17 & \small\code{shipping} & \small\code{shipping_id} & \small\code{order_id} $\to$ \code{orders} (UNIQUE); \code{address_id} $\to$ \code{address} & \small\code{Shipping}\\
\small 18 & \small\code{payment} & \small\code{payment_id} & \small\code{order_id} $\to$ \code{orders} & \small\code{Payment}\\
\small 19 & \small\code{pay_cash} & \small\code{payment_id} & \small\code{payment_id} $\to$ \code{payment} & \small\code{PayCash}\\
\small 20 & \small\code{pay_credit} & \small\code{payment_id} & \small\code{payment_id} $\to$ \code{payment} & \small\code{PayCredit}\\
\bottomrule
\end{longtable}

\subsection{Giải thích quan hệ giữa các bảng (Explain relationships)}

Bảng~\ref{tab:fks} giải thích 21 khóa ngoại. Cột bội số ghi theo thứ tự bảng cha -- bảng con.

\begin{longtable}{@{}P{4.9cm}P{2.4cm}P{1.6cm}P{6.6cm}@{}}
\caption{Khóa ngoại, tên quan hệ và bội số}\label{tab:fks}\\
\toprule
Khóa ngoại (bảng con.cột $\to$ bảng cha) & Tên quan hệ & Bội số & Ý nghĩa\\
\midrule
\endfirsthead
\toprule
Khóa ngoại (bảng con.cột $\to$ bảng cha) & Tên quan hệ & Bội số & Ý nghĩa\\
\midrule
\endhead
\small\code{customer_vip.customer_id} $\to$ \code{customer} & \small is a & \small 1 -- 0..1 & \small Phần riêng của khách VIP\\
\small\code{customer_new.customer_id} $\to$ \code{customer} & \small is a & \small 1 -- 0..1 & \small Phần riêng của khách mới\\
\small\code{address.customer_id} $\to$ \code{customer} & \small has & \small 1 -- 1..* & \small Địa chỉ của khách\\
\small\code{book.product_id} $\to$ \code{product} & \small is a & \small 1 -- 0..1 & \small Phần riêng của sách\\
\small\code{clothes.product_id} $\to$ \code{product} & \small is a & \small 1 -- 0..1 & \small Phần riêng của quần áo\\
\small\code{shoes.product_id} $\to$ \code{product} & \small is a & \small 1 -- 0..1 & \small Phần riêng của giày\\
\small\code{electronics.product_id} $\to$ \code{product} & \small is a & \small 1 -- 0..1 & \small Phần riêng của hàng điện tử\\
\small\code{laptop.product_id} $\to$ \code{electronics} & \small is a & \small 1 -- 0..1 & \small Phần riêng của laptop\\
\small\code{mobile.product_id} $\to$ \code{electronics} & \small is a & \small 1 -- 0..1 & \small Phần riêng của điện thoại\\
\small\code{product_feature.product_id} $\to$ \code{product} & \small indexed by & \small 1 -- 0..1 & \small Vector đặc trưng ảnh của sản phẩm\\
\small\code{cart.customer_id} $\to$ \code{customer} & \small owns & \small 1 -- 0..1 & \small Mỗi khách tối đa một giỏ (UNIQUE)\\
\small\code{cart_item.cart_id} $\to$ \code{cart} & \small contains & \small 1 -- 0..* & \small Dòng của giỏ\\
\small\code{cart_item.product_id} $\to$ \code{product} & \small refers to & \small 1 -- 0..* & \small Sản phẩm của dòng giỏ\\
\small\code{orders.customer_id} $\to$ \code{customer} & \small places & \small 1 -- 0..* & \small Người đặt đơn\\
\small\code{order_item.order_id} $\to$ \code{orders} & \small contains & \small 1 -- 1..* & \small Dòng của đơn\\
\small\code{order_item.product_id} $\to$ \code{product} & \small refers to & \small 1 -- 0..* & \small Sản phẩm đã mua\\
\small\code{shipping.order_id} $\to$ \code{orders} & \small is delivered by & \small 1 -- 0..1 & \small Mỗi đơn tối đa một bản ghi giao (UNIQUE)\\
\small\code{shipping.address_id} $\to$ \code{address} & \small ships to & \small 1 -- 0..* & \small Địa chỉ nhận hàng\\
\small\code{payment.order_id} $\to$ \code{orders} & \small is paid by & \small 1 -- 0..* & \small Các lần thanh toán của đơn\\
\small\code{pay_cash.payment_id} $\to$ \code{payment} & \small is a & \small 1 -- 0..1 & \small Phần riêng của thanh toán tiền mặt\\
\small\code{pay_credit.payment_id} $\to$ \code{payment} & \small is a & \small 1 -- 0..1 & \small Phần riêng của thanh toán thẻ\\
\bottomrule
\end{longtable}

Khóa ngoại và \code{UNIQUE} bảo đảm cận trên của bội số, nhưng không tự bảo đảm cận dưới 1 ở các quan
hệ 1 -- 1..*. Điều kiện ``mỗi đơn có ít nhất một dòng hàng'' và ``mỗi khách có ít nhất một địa chỉ'' được
kiểm tra ở tầng nghiệp vụ: \code{CheckoutService.checkout} từ chối giỏ rỗng (mục 4.4.3). Tương tự, cột
phân loại cho biết bản ghi cha phải có đúng một bản ghi con nào; điều này cũng do tầng nghiệp vụ giữ.

\subsection{Sinh cơ sở dữ liệu (Generate DB)}

Tệp \code{db/ecom_schema.sql} (\sqlLineCount{} dòng, \sqlTableCount{} bảng, cú pháp PostgreSQL) được
sinh bằng \code{tools/gen_sql.py} từ cùng đặc tả bảng đã dùng để dựng Hình~\ref{fig:erd} trong Visual Paradigm. Bảng
được tạo theo thứ tự phụ thuộc khóa ngoại: bảng cha trước, bảng con sau. Đoạn trích dưới đây gồm bảng
\code{product} và hai cặp bảng dòng hàng, thể hiện các quy tắc khóa chính ghép, \code{CHECK} và
\code{ON DELETE CASCADE}.

\codefile{SQL}{../build/excerpts/schema.sql}

Tệp DDL được chạy thử trên SQLite trong bộ nhớ với ràng buộc khóa ngoại bật, nạp dữ liệu mẫu, rồi thử
các câu lệnh vi phạm ràng buộc. Kết quả do script ghi lại như sau:

\txtfile{../build/sql_validation.txt}

Lần chạy này kiểm tra cú pháp và ràng buộc của lược đồ; nó chưa phải phép thử trên máy chủ PostgreSQL
thật.

%% ===================================================================== 4.3
\section{Sơ đồ lớp thiết kế (Class Diagram in Design)}\label{sec:43}

Sơ đồ lớp thiết kế (Hình~\ref{fig:design}) giữ 20 lớp miền của Hình~\ref{fig:analysis} và trả lời câu
hỏi \emph{mỗi lớp giữ dữ liệu gì, với kiểu nào, và chịu trách nhiệm hành vi nào}. Ba enum trạng thái
được thêm vào. Thay đổi chia thành ba phần: thuộc tính, phương thức và quan hệ.

\begin{landscape}
\begin{figure}[p]
\centering
\vpimage{c43_design.png}{0.80\textheight}
\caption[Sơ đồ lớp thiết kế (Design Class Diagram)]{Sơ đồ lớp thiết kế: thuộc tính có kiểu và phạm vi truy cập, phương thức, quan hệ có hướng và tên vai trò.}\label{fig:design}
\end{figure}
\end{landscape}

\subsection{Cập nhật thuộc tính (Update attributes)}

\begin{longtable}{@{}P{3.4cm}P{3.8cm}P{8.7cm}@{}}
\caption{Thuộc tính ở mức phân tích và mức thiết kế}\label{tab:attrs}\\
\toprule
Phân tích (4.1) & Thiết kế (4.3) & Lý do\\
\midrule
\endfirsthead
\toprule
Phân tích (4.1) & Thiết kế (4.3) & Lý do\\
\midrule
\endhead
\small\code{float} cho tiền: \code{price}, \code{amount}, \code{totalAmount}, \code{unitPrice}, \code{cartPrice}, \code{shippingFee}, \code{discountRate}, \code{cashReceived}, \code{changeAmount} & \small\code{BigDecimal} & \small Tiền cần số thập phân chính xác; \code{float} gây sai số làm tròn khi cộng nhiều dòng hàng.\\
\small\code{Date}: \code{dateOfBirth}, \code{registeredDate}, \code{deliveredDate} & \small\code{LocalDate} & \small Chỉ cần ngày, không cần giờ.\\
\small\code{Date}: \code{createdDate}, \code{updatedDate}, \code{orderDate}, \code{paymentDate} & \small\code{LocalDateTime} & \small Cần thời điểm chính xác để sắp xếp và đối soát.\\
\small\code{status: String} của \code{Order}, \code{Payment}, \code{Shipping} & \small Enum \code{OrderStatus}, \code{PaymentStatus}, \code{ShippingStatus} & \small Giới hạn tập giá trị hợp lệ: \code{PLACED}, \code{PROCESSING}, \code{SHIPPED}, \code{DELIVERED}, \code{CANCELLED}; \code{PENDING}, \code{SUCCEEDED}, \code{FAILED}; \code{PREPARING}, \code{IN_TRANSIT}, \code{DELIVERED}. Trình biên dịch phát hiện giá trị sai.\\
\small\code{screenSize: float} & \small\code{double} & \small Kiểu số thực mặc định của Java.\\
\small \code{Address} không có định danh & \small \code{addressId: String} & \small Khi checkout, khách chọn địa chỉ giao bằng mã; \code{Customer.findAddress(addressId)} cần định danh này.\\
\small Phạm vi truy cập chưa xác định & \small Mọi thuộc tính \code{private} (dấu \code{-}) & \small Đóng gói: dữ liệu chỉ được đổi qua phương thức của lớp, nơi kiểm tra quy tắc nghiệp vụ.\\
\bottomrule
\end{longtable}

\subsection{Cập nhật phương thức (Update methods)}

Mỗi lớp miền được giao những hành vi gắn với dữ liệu của nó. Bảng~\ref{tab:ops} liệt kê các phương thức
trên Hình~\ref{fig:design}. Getter và setter không vẽ trên sơ đồ để sơ đồ gọn; chúng được sinh tự động
ở mục 4.5.

\begin{longtable}{@{}P{2.4cm}P{6.6cm}P{6.9cm}@{}}
\caption{Phương thức của các lớp miền}\label{tab:ops}\\
\toprule
Lớp & Phương thức & Trách nhiệm\\
\midrule
\endfirsthead
\toprule
Lớp & Phương thức & Trách nhiệm\\
\midrule
\endhead
\small\code{FullName} & \small\code{+getDisplayName(): String} & \small Ghép họ, tên đệm, tên theo thứ tự tiếng Việt, bỏ phần trống.\\
\small\code{Address} & \small\code{+format(): String} & \small Trả chuỗi địa chỉ để in trên đơn và phiếu giao.\\
\small\code{Customer} & \small\code{+getDisplayName(): String}\newline\code{+getDiscountRate(): BigDecimal}\newline\code{+addAddress(address: Address): void}\newline\code{+findAddress(addressId: String): Address} & \small Họ tên hiển thị; tỷ lệ chiết khấu mặc định bằng 0; thêm địa chỉ; trả địa chỉ thuộc khách hoặc báo lỗi nếu địa chỉ không thuộc khách.\\
\small\code{CustomerVIP} & \small\code{+getDiscountRate(): BigDecimal}\newline\code{+addPoints(points: int): void} & \small Ghi đè để trả \code{discountRate} của khách VIP; cộng điểm, từ chối số âm.\\
\small\code{CustomerNew} & \small\code{+hasWelcomeVoucher(): boolean} & \small Cho biết khách mới còn voucher chào mừng.\\
\small\code{Product} & \small\code{+isAvailable(quantity: int): boolean}\newline\code{+decreaseStock(quantity: int): void}\newline\code|+getCategory(): String {abstract}| & \small Kiểm tra đủ tồn kho; giảm tồn và từ chối khi không đủ; trả nhóm hàng, do từng lớp con cài đặt.\\
\small\code{Book}, \code{Clothes}, \code{Shoes}, \code{Electronics}, \code{Laptop}, \code{Mobile} & \small\code{+getCategory(): String} & \small Trả \code{BOOK}, \code{CLOTHES}, \code{SHOES}, \code{ELECTRONICS}, \code{LAPTOP}, \code{MOBILE}, khớp giá trị của cột \code{product.product_type}.\\
\small\code{Cart} & \small\code{+addItem(product: Product, quantity: int): void}\newline\code{+updateQuantity(productId: String, quantity: int): void}\newline\code{+removeItem(productId: String): void}\newline\code{+calculateTotal(): BigDecimal}\newline\code{+isEmpty(): boolean}\newline\code{+clear(): void} & \small Thêm hàng (cộng dồn nếu sản phẩm đã có trong giỏ); đổi số lượng (số lượng nhỏ hơn hoặc bằng 0 thì xóa dòng); xóa dòng; tính tổng tiền giỏ; kiểm tra giỏ rỗng; làm rỗng giỏ sau checkout.\\
\small\code{CartItem} & \small\code{+CartItem(product: Product, quantity: int)}\newline\code{+getSubtotal(): BigDecimal}\newline\code{+increase(quantity: int): void} & \small Tạo dòng giỏ và lưu \code{cartPrice} bằng giá hiện tại; tính thành tiền; cộng thêm số lượng.\\
\small\code{Order} & \small\code{+addItem(item: OrderItem): void}\newline\code{+addPayment(payment: Payment): void}\newline\code{+calculateTotal(discountRate: BigDecimal): BigDecimal}\newline\code{+cancel(): void}\newline\code{+isOwnedBy(customerId: String): boolean} & \small Thêm dòng đơn; thêm lần thanh toán; tính tổng bằng tổng dòng trừ chiết khấu cộng phí giao, không âm; chỉ hủy được đơn \code{PLACED} hoặc \code{PROCESSING}; kiểm tra đơn có thuộc khách hay không.\\
\small\code{OrderItem} & \small\code|+fromCartItem(item: CartItem, current: Product): OrderItem {static}|\newline\code{+getSubtotal(): BigDecimal} & \small Tạo dòng đơn từ dòng giỏ, chốt \code{unitPrice} bằng giá hiện hành của sản phẩm; tính thành tiền.\\
\small\code{Shipping} & \small\code{+updateStatus(newStatus: ShippingStatus): void}\newline\code{+markDelivered(date: LocalDate): void} & \small Cập nhật trạng thái giao; ghi nhận đã giao và ngày giao.\\
\small\code{Payment} & \small\code|+process(): boolean {abstract}|\newline\code{+markSucceeded(): void}\newline\code{+markFailed(): void} & \small Xử lý thanh toán, do lớp con cài đặt; đặt trạng thái thành công (kèm thời điểm) hoặc thất bại.\\
\small\code{PayCash} & \small\code{+process(): boolean} & \small Đủ tiền thì tính tiền thừa và thành công; thiếu tiền thì thất bại.\\
\small\code{PayCredit} & \small\code{+process(): boolean}\newline\code{+maskedCard(): String} & \small Thành công khi đã có \code{authorizationCode} từ cổng thanh toán; hiển thị thẻ dạng che số.\\
\bottomrule
\end{longtable}

Cách giao phương thức theo nguyên tắc \emph{domain responsibility} của slide A05: hành vi gắn với trạng
thái của một đối tượng thì đặt trong lớp đó. \code{Cart.addItem} và \code{Order.calculateTotal} nằm
trong lớp miền, còn lớp service ở mục 4.4 chỉ điều phối quy trình giữa nhiều đối tượng và kho dữ liệu.
Ba chỗ dùng đa hình:
\begin{itemize}
\item \code{Payment.process()} là phương thức trừu tượng; \code{PayCash} và \code{PayCredit} cài đặt theo
cách riêng, nên service chỉ gọi \code{payment.process()} mà không cần biết hình thức thanh toán.
\item \code{Customer.getDiscountRate()} trả 0, còn \code{CustomerVIP} ghi đè để trả tỷ lệ chiết khấu của
mình. \code{Order.calculateTotal(customer.getDiscountRate())} vì thế đúng với mọi hạng khách.
\item \code{Product.getCategory()} là trừu tượng; mỗi lớp con trả nhóm hàng của mình.
\end{itemize}

\subsection{Cập nhật và giải thích quan hệ (Update relationships -- explain)}

Ở mức phân tích, association chỉ nói hai lớp có liên hệ. Ở mức thiết kế, mỗi association được quyết định
thêm \emph{chiều điều hướng} (lớp nào giữ tham chiếu tới lớp nào) và \emph{tên vai trò} (tên field khi cài
đặt). Bảng~\ref{tab:nav} tóm tắt.

\begin{longtable}{@{}P{3.5cm}P{3.3cm}P{1.9cm}P{6.8cm}@{}}
\caption{Quan hệ ở mức thiết kế}\label{tab:nav}\\
\toprule
Quan hệ & Hướng và vai trò & Bội số & Giải thích\\
\midrule
\endfirsthead
\toprule
Quan hệ & Hướng và vai trò & Bội số & Giải thích\\
\midrule
\endhead
\small\code{Customer} -- \code{FullName} (composition) & \small\code{Customer} $\to$ \code{fullName} & \small 1 -- 1 & \small Khách giữ đối tượng họ tên của mình.\\
\small\code{Customer} -- \code{Address} (aggregation) & \small\code{Customer} $\to$ \code{addresses} & \small 1 -- 1..* & \small Khách giữ danh sách địa chỉ để chọn nơi giao.\\
\small\code{Cart} -- \code{Customer} (owned by) & \small\code{Cart} $\to$ \code{owner} & \small 0..1 -- 1 & \small Đảo chiều so với \emph{owns}: giỏ biết chủ sở hữu để kiểm tra quyền; khách không giữ tham chiếu tới giỏ, giỏ được tìm bằng \code{CartDAO.findByCustomerId}.\\
\small\code{Cart} -- \code{CartItem} (composition) & \small\code{Cart} $\to$ \code{items} & \small 1 -- 0..* & \small Giỏ quản lý các dòng của nó.\\
\small\code{CartItem} -- \code{Product} & \small\code{CartItem} $\to$ \code{product} & \small 0..* -- 1 & \small Dòng giỏ biết sản phẩm; sản phẩm không cần biết nó nằm trong giỏ nào.\\
\small\code{Order} -- \code{Customer} (placed by) & \small\code{Order} $\to$ \code{customer} & \small 0..* -- 1 & \small Đơn biết người đặt (dùng cho \code{isOwnedBy}); danh sách đơn của khách được đọc qua \code{OrderDAO.findByCustomer} thay vì giữ trong \code{Customer}.\\
\small\code{Order} -- \code{OrderItem} (composition) & \small\code{Order} $\to$ \code{items} & \small 1 -- 1..* & \small Đơn quản lý các dòng của nó.\\
\small\code{OrderItem} -- \code{Product} & \small\code{OrderItem} $\to$ \code{product} & \small 0..* -- 1 & \small Dòng đơn biết sản phẩm đã mua.\\
\small\code{Order} -- \code{Shipping} (delivered by) & \small\code{Order} $\to$ \code{shipping} & \small 1 -- 0..1 & \small Đơn giữ thông tin giao hàng.\\
\small\code{Shipping} -- \code{Address} (ships to) & \small\code{Shipping} $\to$ \code{address} & \small 0..* -- 1 & \small Lần giao biết địa chỉ đích.\\
\small\code{Order} -- \code{Payment} (paid by) & \small\code{Order} $\to$ \code{payments} & \small 1 -- 0..* & \small Đơn giữ các lần thanh toán.\\
\small Kế thừa & \small \code{extends} & & \small Giữ nguyên sáu nhánh kế thừa; \code{Product} và \code{Payment} là lớp trừu tượng.\\
\small\code{OrderItem} $\dashrightarrow$ \code{CartItem} & \small «use» & & \small \code{OrderItem.fromCartItem} đọc dòng giỏ để tạo dòng đơn.\\
\bottomrule
\end{longtable}

Điều hướng một chiều giảm phụ thuộc vòng: nếu cả \code{Customer} và \code{Order} giữ tham chiếu tới nhau,
mỗi lần đọc khách sẽ kéo theo toàn bộ đơn. Phụ thuộc «create» \emph{checkout} \code{Cart}
$\dashrightarrow$ \code{Order} của mức phân tích không còn nằm giữa hai lớp miền. Ở mức thiết kế, việc
tạo đơn từ giỏ được tách thành hai phần: \code{OrderItem.fromCartItem} chuyển từng dòng
(phụ thuộc «use» \code{OrderItem} $\dashrightarrow$ \code{CartItem}), còn \code{CheckoutService} điều
phối toàn bộ quy trình và tạo \code{Order} (phụ thuộc «create» ở Hình~\ref{fig:layerclasses}).

%% ===================================================================== 4.4
\section{Tầng, lớp và quan hệ (Layer + Class + Relationship)}\label{sec:44}

\subsection{Bốn tầng và quy tắc phụ thuộc (Layers)}

\begin{figure}[H]
\centering
\vpimage{c44_layers.png}{0.42\textheight}
\caption[Bốn tầng và package của hệ thống]{Bốn layer (package) của hệ thống và phụ thuộc giữa chúng.}\label{fig:layers}
\end{figure}

Theo slide A05, hệ thống được chia thành bốn tầng Presentation, Application, Domain và Infrastructure
(Hình~\ref{fig:layers}). Mũi tên nét đứt là phụ thuộc: Presentation dùng Application; Application dùng
Domain và Infrastructure; Infrastructure dùng Domain để đọc và ghi đối tượng miền. Không có mũi tên từ
Presentation xuống Infrastructure, nghĩa là giao diện không truy cập cơ sở dữ liệu trực tiếp. Domain
không phụ thuộc tầng nào, nên quy tắc nghiệp vụ không bị buộc vào framework hay cơ sở dữ liệu. Bên trong
mỗi layer là các package con trùng với package Java ở mục 4.5: \code{controller}; \code{service},
\code{search}; \code{customer}, \code{product}, \code{cart}, \code{order}, \code{payment}; \code{dao},
\code{dao.jdbc}, \code{search}.

\begin{longtable}{@{}P{2.1cm}P{3.5cm}P{3.9cm}P{3.1cm}P{2.5cm}@{}}
\caption{Trách nhiệm của từng tầng}\label{tab:layers}\\
\toprule
Tầng & Package Java & Lớp & Trách nhiệm & Không được làm\\
\midrule
\endfirsthead
\toprule
Tầng & Package Java & Lớp & Trách nhiệm & Không được làm\\
\midrule
\endhead
\small Presentation & \small\code{com.ptit.ecom.presentation.controller} & \small\code{CartController}, \code{CheckoutController}, \code{OrderController}, \code{SearchController} & \small Nhận yêu cầu, lấy \code{customerId} từ phiên đăng nhập, gọi service, trả kết quả & \small Truy cập CSDL; chứa quy tắc nghiệp vụ\\
\small Application & \small\code{com.ptit.ecom.application.service}, \code{com.ptit.ecom.application.search} & \small\code{CartService}, \code{CheckoutService}, \code{OrderService}; \code{SearchService}, \code{QueryService}, \code{RankingService}, \code{CommonQuery}, \code{SearchResult} & \small Điều phối use case, xác định ranh giới transaction, gọi DAO & \small Viết SQL; giữ trạng thái giao diện\\
\small Domain & \small\code{com.ptit.ecom.domain.customer}, \code{.product}, \code{.cart}, \code{.order}, \code{.payment} & \small 20 lớp miền của mục 4.3 và ba enum trạng thái & \small Dữ liệu và quy tắc nghiệp vụ (\code{addItem}, \code{calculateTotal}, \code{cancel}, \code{process}) & \small Phụ thuộc framework, CSDL hay giao diện\\
\small Infrastructure & \small\code{com.ptit.ecom.infrastructure.dao}, \code{.dao.jdbc}, \code{.search} & \small\code{CustomerDAO}, \code{ProductDAO}, \code{CartDAO}, \code{OrderDAO}; \code{JdbcCustomerDAO}, \code{JdbcProductDAO}, \code{JdbcCartDAO}, \code{JdbcOrderDAO}; \code{VectorIndex} & \small Lưu trữ, câu lệnh SQL, adapter tìm ảnh & \small Quyết định giá, quyền hay trạng thái đơn\\
\bottomrule
\end{longtable}

Package không đồng nghĩa với layer (slide A05): layer là trách nhiệm kiến trúc, package là cách tổ chức
mã. Ở đây mỗi layer là một package gốc, bên trong chia tiếp theo chức năng (\code{service} và
\code{search}) hoặc theo nhóm miền (\code{cart}, \code{order}, \code{product}, \code{customer},
\code{payment}). Controller trả về đối tượng miền (\code{Cart}, \code{Order}) cho phía hiển thị, tức
Presentation có đọc kiểu của Domain; bản hoàn chỉnh có thể thay bằng DTO nếu cần ẩn bớt dữ liệu.

\subsection{Lớp và quan hệ trong từng tầng (Classes and relationships)}

\begin{landscape}
\begin{figure}[p]
\centering
\vpimage{c44_classes.png}{0.80\textheight}
\caption[Lớp và quan hệ trong từng tầng]{Lớp trong từng tầng và quan hệ \emph{use}, \emph{create}, \emph{realization} giữa chúng.}\label{fig:layerclasses}
\end{figure}
\end{landscape}

Hình~\ref{fig:layerclasses} xếp 44 lớp theo bốn tầng: Presentation ở trên, Application ở giữa,
Infrastructure ở dưới và Domain bên phải. Mỗi nhãn ghi tên tầng hoặc package Java của nhóm lớp ngay dưới
nó; trong project VP, mỗi lớp nằm đúng package con của Hình~\ref{fig:layers}. Lớp được vẽ thu gọn (chỉ
tên); thuộc tính và phương thức đầy đủ nằm ở Hình~\ref{fig:design}. Có ba loại quan hệ.

\textbf{Phụ thuộc «use» giữa các tầng.} Mỗi controller dùng đúng một service:
\code{CartController} $\to$ \code{CartService}, \code{CheckoutController} $\to$ \code{CheckoutService},
\code{OrderController} $\to$ \code{OrderService}, \code{SearchController} $\to$ \code{SearchService}.
Service dùng các interface DAO: \code{CartService} dùng \code{CartDAO} và \code{ProductDAO};
\code{CheckoutService} dùng \code{CartDAO}, \code{OrderDAO}, \code{ProductDAO}, \code{CustomerDAO};
\code{OrderService} dùng \code{OrderDAO}; \code{SearchService} dùng \code{QueryService},
\code{RankingService}, \code{ProductDAO} và \code{VectorIndex}. \code{CartService} còn dùng \code{Cart}
để gọi hành vi miền. Mỗi phụ thuộc được cài đặt bằng một field nhận qua constructor, ví dụ
\code{CartService(cartDAO: CartDAO, productDAO: ProductDAO)}.

\textbf{Phụ thuộc «create».} \code{CheckoutService} $\dashrightarrow$ \code{Order}: service tạo đơn mới.
Đây là bản thiết kế của phép biến đổi checkout ở mức phân tích.

\textbf{Realization.} \code{JdbcCustomerDAO}, \code{JdbcProductDAO}, \code{JdbcCartDAO} và
\code{JdbcOrderDAO} cài đặt bốn interface tương ứng. Service chỉ phụ thuộc interface, không phụ thuộc lớp
JDBC (nguyên tắc dependency inversion). Nhờ vậy có thể đổi hệ quản trị CSDL, hoặc thay bằng DAO trong bộ
nhớ khi kiểm thử service, mà không sửa tầng Application.

\subsection{Checkout đi qua các tầng (Explain)}

Phương thức \code{CheckoutService.checkout} là use case trung tâm của slide A05. Trình tự dưới đây bám
theo thân phương thức trong đặc tả thiết kế.
\begin{enumerate}
\item Khách bấm \emph{Đặt hàng}. \code{CheckoutController.checkout(customerId, addressId, paymentMethod)}
nhận yêu cầu; \code{customerId} lấy từ phiên đăng nhập.
\item Controller tạo \code{PayCash} hoặc \code{PayCredit} theo \code{paymentMethod}, rồi gọi
\code{CheckoutService.checkout(customerId, addressId, payment)}.
\item Service đọc giỏ bằng \code{CartDAO.findByCustomerId(customerId)}; giỏ không có hoặc rỗng thì từ chối.
\item Service đọc khách bằng \code{CustomerDAO.findById(customerId)}.
\item Service tạo \code{Order} mới với trạng thái \code{PLACED} và thời điểm đặt hiện tại.
\item Với mỗi \code{CartItem}, service đọc lại sản phẩm bằng \code{ProductDAO.findById}, kiểm tra
\code{isAvailable(quantity)}, rồi thêm \code{OrderItem.fromCartItem(item, current)}. Giá được chốt theo
giá hiện hành của catalog, không lấy giá trong giỏ.
\item Service tạo \code{Shipping} với địa chỉ \code{customer.findAddress(addressId)} (địa chỉ phải thuộc
khách), trạng thái \code{PREPARING}, rồi gắn vào đơn.
\item \code{order.calculateTotal(customer.getDiscountRate())} tính tổng tiền theo hạng khách.
\item Service gán số tiền cho \code{payment}, đặt trạng thái \code{PENDING} và gọi
\code{order.addPayment(payment)}.
\item \code{OrderDAO.save(order)}, \code{ProductDAO.decreaseStock} cho từng dòng và
\code{CartDAO.clear(cartId)} chạy trong một transaction; service trả \code{Order} cho controller.
\end{enumerate}

\IfFileExists{../figures/vp/c44_checkout_seq.png}{%
\begin{figure}[H]
\centering
\vpimage{c44_checkout_seq.png}{0.60\textheight}
\caption[Sequence diagram thiết kế cho checkout]{Sequence diagram thiết kế: checkout đi qua Controller, Service, Domain và DAO.}\label{fig:seq}
\end{figure}
Hình~\ref{fig:seq} vẽ lại trình tự trên bằng sequence diagram thiết kế trong Visual Paradigm. Mỗi lifeline
có kiểu là một lớp của Hình~\ref{fig:layerclasses}, nên mỗi thông điệp là một phương thức đã khai báo ở
mục 4.3 hoặc 4.4: thông điệp 2 là \code{CheckoutService.checkout}, các thông điệp 3, 5, 8, 12, 13, 14 đi
xuống interface DAO, còn 7, 10, 11 gọi hành vi của lớp miền \code{Order}. Khung \emph{loop} lặp bước đọc
lại sản phẩm và tạo \code{OrderItem} cho từng dòng giỏ. Đường trả về nét đứt mang kết quả ngược lên tầng
trên; không có thông điệp nào đi từ Controller thẳng xuống DAO.}{}

\textbf{Vì sao checkout nhận customerId thay vì cartID.} Slide A05 gợi ý
\code{CheckoutService.checkout(cartID): Order}. Thiết kế này dùng \code{customerId} lấy từ phiên đăng
nhập. Mỗi khách có tối đa một giỏ (Customer -- Cart là 1 -- 0..1), nên \code{customerId} vừa xác định
được giỏ, vừa bảo đảm khách chỉ checkout giỏ của chính mình, thay vì tin một \code{cartId} do trình
duyệt gửi lên. \code{addressId} chọn địa chỉ giao trong danh sách địa chỉ của khách; \code{payment} là
hình thức thanh toán khách đã chọn.

\subsection{Gói tìm kiếm search (Explain)}

Gói \code{com.ptit.ecom.application.search} đưa phân hệ tìm kiếm đa phương thức của prototype A6 vào
kiến trúc chung. \code{SearchController.search(mode, text, image)} chọn một trong ba đường:
\begin{itemize}
\item \textbf{Văn bản:} \code{SearchService.searchText} gọi \code{QueryService.buildTextQuery} để chuẩn
hóa truy vấn thành \code{CommonQuery}, lấy ứng viên bằng \code{ProductDAO.searchByKeyword}, chấm điểm
khớp tên, rồi gọi \code{RankingService.rank}.
\item \textbf{Giọng nói:} \code{searchVoice} nhận transcript và đi chung đường văn bản, đúng như cách
prototype A6 xử lý transcript.
\item \textbf{Ảnh:} \code{searchImage} gọi \code{VectorIndex.embed} để lấy vector đặc trưng, lấy các sản
phẩm gần nhất bằng \code{VectorIndex.nearest}, đọc lại sản phẩm từ \code{ProductDAO}, rồi xếp hạng.
\end{itemize}
Bước lấy ứng viên (retrieval) và bước xếp hạng (ranking) được tách thành hai lớp. \code{RankingService}
loại sản phẩm hết hàng, sắp xếp theo điểm và lấy năm kết quả đầu, nên quy tắc tồn kho vẫn do dữ liệu
catalog quyết định. \code{VectorIndex} là adapter tới engine đặc trưng ảnh (Apple Vision trong
prototype A6); vector được lưu ở bảng \code{product_feature} kèm \code{model_version}.

%% ===================================================================== 4.5
\section{Sinh mã nguồn từ mô hình 4.4 (Generate code from 4.4)}\label{sec:45}

\subsection{Cách sinh mã từ mô hình 4.4 (Generate code)}

Mã Java được sinh bằng script \code{tools/gen_java.py}. Script đọc cùng tệp đặc tả \code{tools/model.py}
đã dùng để dựng các sơ đồ 4.3 và 4.4 trong Visual Paradigm, nên package, lớp, field và chữ ký phương thức
trùng với sơ đồ. Bài không dùng nút Generate Code của Visual Paradigm; nếu phiên bản VP có chức năng sinh
mã (Tools $>$ Code), có thể sinh lại từ project để đối chiếu. Kết quả gồm \javaFileCount{} tệp Java
(\javaLineCount{} dòng) trong \javaPackageCount{} package, đặt tại \code{code/src}.

\begin{longtable}{@{}P{4.5cm}P{5.4cm}P{6.0cm}@{}}
\caption{Quy tắc ánh xạ từ mô hình UML sang Java}\label{tab:uml2java}\\
\toprule
Phần tử UML (mục 4.3, 4.4) & Mã Java & Ví dụ\\
\midrule
\endfirsthead
\toprule
Phần tử UML (mục 4.3, 4.4) & Mã Java & Ví dụ\\
\midrule
\endhead
\small Package của layer & \small Package \code{com.ptit.ecom.<layer>.<nhóm>} & \small\code{com.ptit.ecom.domain.cart}\\
\small Lớp / lớp trừu tượng & \small\code{public class} / \code{public abstract class} & \small\code{Product}, \code{Payment} là abstract\\
\small Lớp «interface» & \small\code{public interface} & \small\code{OrderDAO}\\
\small Enum & \small\code{public enum} & \small\code{OrderStatus}\\
\small Thuộc tính \code{-name: Type} & \small Field \code{private Type name;} & \small\code{private BigDecimal unitPrice;}\\
\small Đầu association có vai trò, bội số 1 hoặc 0..1 & \small Field kiểu lớp đích & \small\code{private Customer owner;}\\
\small Đầu association bội số 0..* hoặc 1..* & \small Field \code{List} khởi tạo rỗng & \small\code{private List<CartItem> items}\newline\code{= new ArrayList<>();}\\
\small Generalization & \small\code{extends} & \small\code{CustomerVIP extends Customer}\\
\small Realization & \small\code{implements} & \small\code{JdbcOrderDAO implements OrderDAO}\\
\small Phương thức & \small Phương thức \code{public}/\code{private}, cùng chữ ký; \code|{abstract}| và \code|{static}| giữ nguyên & \small\code{public static OrderItem fromCartItem(...)}\\
\small Phương thức trùng tên lớp & \small Constructor & \small\code{CartItem(Product, int)}\\
\small Phương thức chưa có thân trong đặc tả & \small Thân tạm {\footnotesize\code{throw new}\newline\code{UnsupportedOperationException(}\newline\code{"TODO: ...")}} & \small\code{JdbcOrderDAO.findByCustomer}\\
\small Field (không vẽ accessor trên sơ đồ) & \small Getter và setter sinh tự động & \small\code{getItems()}, \code{setStatus(...)}\\
\bottomrule
\end{longtable}

\subsection{Giải thích các package (Explain packages)}

Cây thư mục mã nguồn do script sinh ra:

\txtfile{../build/code_tree.txt}

\begin{longtable}{@{}P{5.5cm}P{2.4cm}P{8.0cm}@{}}
\caption{Nội dung từng package}\label{tab:packages}\\
\toprule
Package (\code{com.ptit.ecom.…}) & Tầng & Nội dung\\
\midrule
\endfirsthead
\toprule
Package (\code{com.ptit.ecom.…}) & Tầng & Nội dung\\
\midrule
\endhead
\small\code{presentation.controller} & \small Presentation & \small Bốn controller nhận yêu cầu của giỏ, checkout, đơn và tìm kiếm.\\
\small\code{application.service} & \small Application & \small \code{CartService}, \code{CheckoutService}, \code{OrderService}: điều phối use case giỏ, đặt hàng, tra cứu và hủy đơn.\\
\small\code{application.search} & \small Application & \small \code{SearchService}, \code{QueryService}, \code{RankingService} và hai lớp dữ liệu \code{CommonQuery}, \code{SearchResult}.\\
\small\code{domain.customer} & \small Domain & \small \code{Customer}, \code{FullName}, \code{Address}, \code{CustomerVIP}, \code{CustomerNew}.\\
\small\code{domain.product} & \small Domain & \small \code{Product} và sáu lớp con.\\
\small\code{domain.cart} & \small Domain & \small \code{Cart}, \code{CartItem}.\\
\small\code{domain.order} & \small Domain & \small \code{Order}, \code{OrderItem}, \code{Shipping} và hai enum \code{OrderStatus}, \code{ShippingStatus}.\\
\small\code{domain.payment} & \small Domain & \small \code{Payment}, \code{PayCash}, \code{PayCredit}, enum \code{PaymentStatus}.\\
\small\code{infrastructure.dao} & \small Infrastructure & \small Bốn interface DAO mà tầng Application phụ thuộc.\\
\small\code{infrastructure.dao.jdbc} & \small Infrastructure & \small Bốn lớp JDBC cài đặt các interface trên.\\
\small\code{infrastructure.search} & \small Infrastructure & \small \code{VectorIndex}, adapter tới engine đặc trưng ảnh.\\
\bottomrule
\end{longtable}

Hướng \code{import} giữa các package đi theo Hình~\ref{fig:layers}: \code{presentation} import
\code{application}; \code{application} import \code{domain} và \code{infrastructure.dao};
\code{infrastructure} import \code{domain}; các package \code{domain} chỉ import lẫn nhau và thư viện
chuẩn của Java.

\subsection{Giải thích mã của các lớp (Class code)}

Các đoạn dưới đây được trích từ mã đã sinh; phần accessor và các thành phần khác được lược bớt.

\textbf{Lớp miền \texttt{Cart}.}

\codefile{Java}{../build/excerpts/cart.java}

Field \code{owner} và \code{items} được sinh từ hai đầu association \code{Cart} $\to$ \code{Customer} và
composition \code{Cart} $\to$ \code{CartItem}. \code{addItem} cộng dồn số lượng nếu sản phẩm đã có
trong giỏ, nên mỗi sản phẩm chỉ có một dòng giỏ; điều này khớp với khóa chính ghép
\code{(cart_id, product_id)} của bảng \code{cart_item}. \code{calculateTotal} cộng tiền bằng
\code{BigDecimal}.

\textbf{Phương thức tĩnh \texttt{OrderItem.fromCartItem}.}

\codefile{Java}{../build/excerpts/orderitem.java}

Dòng đơn lấy số lượng từ dòng giỏ nhưng lấy giá từ đối tượng \code{current} vừa đọc lại từ catalog.
\code{unitPrice} vì thế là giá chốt tại thời điểm đặt, đúng ý nghĩa đã phân tích ở mục 4.1.

\textbf{Service \texttt{CheckoutService}.}

\codefile{Java}{../build/excerpts/checkout.java}

Bốn field DAO là các phụ thuộc «use» của Hình~\ref{fig:layerclasses}. Thân phương thức là mười bước của
mục 4.4.3: service chỉ điều phối, còn việc kiểm tra tồn kho, tạo dòng đơn, tính tổng và quản lý thanh
toán được giao cho \code{Product}, \code{OrderItem}, \code{Order} và \code{Payment}.

\textbf{DAO \texttt{OrderDAO} và \texttt{JdbcOrderDAO}.}

\codefile{Java}{../build/excerpts/orderdao.java}

Interface \code{OrderDAO} thuộc \code{infrastructure.dao}; lớp \code{JdbcOrderDAO} cài đặt nó
(realization). \code{findByIdAndCustomer} dùng \code{PreparedStatement} với tham số, nên giá trị đầu vào
không bị ghép thẳng vào câu SQL. Điều kiện \code{customer_id = ?} chỉ trả đơn thuộc khách đang đăng nhập,
cài đặt yêu cầu FR-ORDER-02.

\textbf{Lớp trừu tượng \texttt{Payment} và lớp con \texttt{PayCash}.}

\codefile{Java}{../build/excerpts/payment.java}

\code{Payment.process()} là phương thức trừu tượng; \code{PayCash} ghi đè để kiểm tra số tiền nhận và
tính tiền thừa. Đây là đa hình đã nêu ở mục 4.3.2.

\textbf{Gói search.}

\codefile{Java}{../build/excerpts/search.java}

\code{searchImage} cho thấy phần retrieval (\code{VectorIndex.nearest} và \code{ProductDAO.findById})
tách khỏi phần ranking (\code{RankingService.rank}). \code{SearchController.search} chỉ chọn đường xử lý
theo \code{mode}.

\subsection{Biên dịch và kiểm tra nhanh}

Mã sinh ra được biên dịch bằng \code{javac}. Một chương trình kiểm tra nhanh (\code{code/test/DomainSmokeTest.java})
gọi các phương thức miền có thân cài đặt với dữ liệu mẫu. Kết quả do chương trình in ra:

\txtfile{../build/smoke_test.txt}

Các kiểm tra này chỉ bao phủ lớp miền. Các phương thức JDBC chưa được chạy với cơ sở dữ liệu.

\subsection{Giới hạn của mã sinh ra}

\begin{itemize}
\item Đây là khung mã từ mô hình thiết kế. Các phương thức chưa có thân trong đặc tả, gồm phần lớn phương
thức của \code{JdbcCustomerDAO}, \code{JdbcProductDAO}, \code{JdbcCartDAO}, \code{JdbcOrderDAO} và hai
phương thức của \code{VectorIndex}, còn ném \code{UnsupportedOperationException("TODO: ...")}.
\item \code{CheckoutService.checkout} cần một transaction bao ba thao tác lưu đơn, giảm tồn và làm rỗng
giỏ; hiện mỗi DAO tự mở kết nối riêng, nên tầng Infrastructure còn phải cung cấp cơ chế transaction.
\item Phí giao hàng đang đặt bằng 0; cổng thanh toán thẻ và giao diện người dùng chưa được cài đặt.
\end{itemize}

Việc tiếp theo cho buổi đánh giá code ngày 05/11: cài đặt các DAO trên lược đồ ở mục 4.2 (hoặc DAO
trong bộ nhớ cho demo), thêm transaction cho checkout, viết unit test cho \code{Cart}, \code{Order},
\code{Payment}, và nối \code{VectorIndex} với engine ảnh của prototype A6.

%% ===================================================================== 4.6
\section{Kết luận chương (Conclusion)}\label{sec:46}

Chương 4 đi một mạch từ mô hình nghiệp vụ tới mã nguồn. Sơ đồ phân tích 20 lớp (4.1) được chuyển thành
data model và 20 bảng có khóa chính, khóa ngoại và ràng buộc (4.2). Sơ đồ lớp thiết kế bổ sung kiểu dữ
liệu, phương thức và chiều điều hướng (4.3). Kiến trúc bốn tầng đặt lớp miền cạnh controller, service và
DAO (4.4). Cuối cùng, mã Java được sinh từ cùng mô hình đó (4.5). Vì sơ đồ, SQL và mã cùng đọc một tệp
đặc tả, tên lớp, tên bảng và chữ ký phương thức nhất quán giữa các mục.

\begin{longtable}{@{}P{2.9cm}P{6.3cm}P{3.5cm}P{2.8cm}@{}}
\caption{Truy vết từ yêu cầu tới thiết kế, dữ liệu và mã}\label{tab:trace}\\
\toprule
Yêu cầu & Lớp và phương thức & Bảng, cột & Mục\\
\midrule
\endfirsthead
\toprule
Yêu cầu & Lớp và phương thức & Bảng, cột & Mục\\
\midrule
\endhead
\small FR-CART-01: thêm sản phẩm còn bán vào giỏ & \small\code{CartController.addItem} $\to$ \code{CartService.addToCart} $\to$ \code{Cart.addItem}, \code{Product.isAvailable} & \small\code{cart}, \code{cart_item} (\code{quantity > 0}) & \small 4.3, 4.4, 4.5\\
\small FR-ORDER-01: tạo đơn từ giỏ hợp lệ, chốt giá & \small\code{CheckoutService.checkout}, \code{OrderItem.fromCartItem}, \code{Order.calculateTotal} & \small\code{orders}, \code{order_item.unit_price} & \small 4.3, 4.4.3, 4.5\\
\small FR-ORDER-02: khách chỉ xem đơn của mình & \small\code{OrderService.getOrder} $\to$ \code{OrderDAO.findByIdAndCustomer}; \code{Order.isOwnedBy} & \small\code{orders.customer_id} & \small 4.4, 4.5\\
\small FR-SEARCH-01: tìm theo từ khóa & \small\code{SearchService.searchText}, \code{QueryService.buildTextQuery}, \code{RankingService.rank} & \small\code{product} & \small 4.4.4, 4.5\\
\small FR-SEARCH-03: tìm bằng ảnh & \small\code{SearchService.searchImage}, \code{VectorIndex.nearest} & \small\code{product_feature} & \small 4.2, 4.4.4\\
\small D-01: không lưu số thẻ đầy đủ & \small\code{PayCredit.cardLast4}, \code{authorizationCode}, \code{maskedCard} & \small\code{pay_credit.card_last4} & \small 4.1, 4.2, 4.3\\
\bottomrule
\end{longtable}

\textbf{Hạn chế.} Mô hình dữ liệu chưa có lịch sử chuyển trạng thái đơn và chưa có tồn kho theo kho
hàng. Phí giao hàng chưa được thiết kế. Ràng buộc ``đúng một bản ghi con'' của cột phân loại và cận dưới
1..* vẫn do tầng nghiệp vụ giữ. Mã ở mục 4.5 là khung có một phần thân cài đặt, chưa phải ứng dụng chạy
được.

\textbf{Kế hoạch.} Ngày 29/10 đánh giá phần thiết kế; ngày 05/11 đánh giá code, với các việc nêu ở cuối
mục 4.5; ngày 12/11 và 19/11 đánh giá tiểu luận final và demo của từng sinh viên. Khi ghép vào tiểu luận
final, chương này thay chương Design của bản V1.

%% ================================================================ appendix
\section*{Phụ lục. Tệp nộp kèm và cách tái tạo}
\addcontentsline{toc}{section}{Phụ lục. Tệp nộp kèm và cách tái tạo}

\begin{table}[H]
\centering\small
\begin{tabular}{@{}P{5.6cm}P{10.0cm}@{}}
\toprule
Tệp, thư mục (trong \code{tieuluan/}) & Nội dung\\
\midrule
\code{vp/tieuluan_design.vpp} & Project Visual Paradigm chứa toàn bộ sơ đồ của chương, mở và chỉnh sửa được\\
\code{figures/vp/*.png}, \code{*.svg} & Ảnh sơ đồ xuất từ project\\
\code{db/ecom_schema.sql} & DDL PostgreSQL của 20 bảng\\
\code{code/src} & Mã Java sinh từ mô hình 4.4\\
\code{tools/model.py} & Đặc tả chung của lớp, quan hệ, bảng\\
\code{chapter4/design_04.tex} & Mã nguồn LaTeX của chương\\
\bottomrule
\end{tabular}
\end{table}

Lệnh tái tạo, chạy trong thư mục \code{tieuluan/}: \code{tools/build_vp.py} dựng project VP và xuất
hình; \code{tools/gen_sql.py} sinh \code{db/ecom_schema.sql} và kiểm tra trên SQLite;
\code{tools/gen_java.py} sinh \code{code/src}, biên dịch và chạy kiểm tra nhanh; lệnh cuối biên dịch
chương này.
\begin{Verbatim}[fontsize=\small,frame=single]
python3 tools/build_vp.py
python3 tools/gen_sql.py
python3 tools/gen_java.py
tectonic -X compile chapter4/design_04.tex --outdir build
\end{Verbatim}

\end{document}
